\documentclass[../DoAn.tex]{subfiles}
\begin{document}

\section{Cấu hình WSL2 Mirrored Networking}
\label{sec:pla-wsl}

WSL2 Mirrored Networking cho phép WSL2 interface nhận traffic trực tiếp từ mạng LAN vật lý, thay vì qua NAT của Hyper-V. Đây là yêu cầu bắt buộc để CICFlowMeter thu thập flow features giống với môi trường mạng thực.

\textbf{Kích hoạt Mirrored Networking:}

Tạo hoặc chỉnh sửa file \texttt{\%USERPROFILE\%\textbackslash.wslconfig} trên Windows 11:

\begin{lstlisting}[language=bash, caption={Cấu hình .wslconfig để bật Mirrored Networking}]
[wsl2]
networkingMode=mirrored
\end{lstlisting}

Sau đó khởi động lại WSL2:

\begin{lstlisting}[language=bash, caption={Khởi động lại WSL2}]
wsl --shutdown
wsl
\end{lstlisting}

\textbf{Kiểm tra cấu hình:}

\begin{lstlisting}[language=bash, caption={Kiểm tra IP LAN trên WSL2}]
# Trong WSL2 Ubuntu -- kiem tra interface eth0 nhan IP LAN
ip addr show eth0
# Ky vong: IP thuoc dai 192.168.x.x (LAN), khong phai 172.x.x.x (NAT)
\end{lstlisting}

\textbf{Lưu ý:} Mirrored Networking yêu cầu Windows 11 Build 22621.2359 trở lên và WSL 2.0.0 trở lên. Kiểm tra phiên bản WSL bằng lệnh \texttt{wsl --version}.

\section{Cài đặt và cấu hình Snort 3}
\label{sec:pla-snort}

\textbf{Cài đặt dependencies và Snort 3:}

\begin{lstlisting}[language=bash, caption={Cài đặt Snort 3 từ source trên Ubuntu 22.04}]
sudo apt update
sudo apt install -y build-essential libpcap-dev libpcre2-dev \
    libdnet-dev zlib1g-dev liblzma-dev openssl libssl-dev \
    cmake flex bison git pkg-config

# LibDAQ
git clone https://github.com/snort3/libdaq.git
cd libdaq && ./bootstrap && ./configure && make && sudo make install
cd ..

# Snort 3
git clone https://github.com/snort3/snort3.git
cd snort3
./configure_cmake.sh --prefix=/usr/local --enable-tcmalloc
cd build && make -j$(nproc) && sudo make install
\end{lstlisting}

\textbf{Cấu hình snort.lua cho Testbed:}

\begin{lstlisting}[language={[5.1]Lua}, caption={Cấu hình tối thiểu snort.lua cho môi trường Testbed}]
-- /usr/local/etc/snort/snort.lua
HOME_NET = '192.168.1.0/24'
EXTERNAL_NET = 'any'

ips = {
    enable_builtin_rules = true,
    include = RULE_PATH .. '/local.rules',
}

alert_fast = {
    file = true,
    packet = false,
}
\end{lstlisting}

\textbf{Rules tùy chỉnh (local.rules):}

\begin{lstlisting}[language=bash, caption={Rules Snort phát hiện DoS và BruteForce SSH}]
# Phat hien DoS HTTP flood
alert tcp $EXTERNAL_NET any -> $HOME_NET 80 \
    (msg:"DoS HTTP Flood detected"; \
     flow:to_server; \
     threshold:type threshold,track by_src,count 100,seconds 1; \
     sid:1000001; rev:1;)

# Phat hien SSH BruteForce
alert tcp $EXTERNAL_NET any -> $HOME_NET 22 \
    (msg:"SSH BruteForce attempt"; \
     flow:to_server,established; \
     threshold:type threshold,track by_src,count 10,seconds 5; \
     sid:1000002; rev:1;)
\end{lstlisting}

\textbf{Chạy Snort trên interface WSL:}

\begin{lstlisting}[language=bash, caption={Khởi động Snort 3 trên interface eth0}]
sudo snort -c /usr/local/etc/snort/snort.lua \
    -i eth0 \
    -l /var/log/snort \
    --daq-dir /usr/local/lib/daq \
    -A fast &
\end{lstlisting}

\section{Triển khai CICFlowMeter trên WSL}
\label{sec:pla-cicflowmeter}

CICFlowMeter trích xuất 77--80 đặc trưng thống kê flow từ PCAP file hoặc live interface.

\textbf{Cài đặt (yêu cầu Java 8):}

\begin{lstlisting}[language=bash, caption={Cài đặt CICFlowMeter trên Ubuntu}]
sudo apt install openjdk-8-jdk

# Dung ban fork voi bug fixes cho Linux
git clone https://github.com/datthinh1801/cicflowmeter.git
cd cicflowmeter
mvn package -DskipTests
\end{lstlisting}

\textbf{Thu thập flow từ PCAP:}

\begin{lstlisting}[language=bash, caption={Chạy CICFlowMeter trên file PCAP}]
java -jar target/CICFlowMeter-4.0-SNAPSHOT.jar \
    /path/to/capture.pcap \
    /path/to/output/flows/
\end{lstlisting}

\textbf{Xử lý đặc trưng trong Python:}

\begin{lstlisting}[language=Python, caption={Load flow CSV và apply PowerTransformer V8.5}]
import pandas as pd
from sklearn.preprocessing import PowerTransformer
import joblib

# Load scaler da train tren Testbed V8.5
scaler = joblib.load('models/v8_5_scaler.pkl')

# Load flow CSV tu CICFlowMeter
flows = pd.read_csv('/tmp/flows/capture_ICTAI_Flow.csv')

# Xu ly Inf/NaN
flows.replace([float('inf'), float('-inf')], float('nan'), inplace=True)
flows.dropna(inplace=True)

# Transform va phan loai
X = scaler.transform(flows[FEATURE_COLUMNS])
\end{lstlisting}

\section{Hướng dẫn khởi chạy pipeline End-to-End}
\label{sec:pla-pipeline}

\begin{lstlisting}[language=bash, caption={Script khởi động pipeline Hybrid IDS hoàn chỉnh}]
# 1. Khoi dong Snort (background)
sudo snort -c /usr/local/etc/snort/snort.lua -i eth0 \
    -l /var/log/snort -A fast &

# 2. Khoi dong capture PCAP (background)
sudo tcpdump -i eth0 -w /tmp/live_capture.pcap &

# 3. Khoi dong inference daemon (Python)
python3 testbed_inference.py \
    --pcap /tmp/live_capture.pcap \
    --model models/v8_5_model.pt \
    --scaler models/v8_5_scaler.pkl \
    --output /var/log/ids/ml_alerts.log &

echo "IDS Pipeline started. Snort + FTT V8.5 active."
\end{lstlisting}

\textbf{Dừng pipeline:}

\begin{lstlisting}[language=bash, caption={Dừng tất cả tiến trình pipeline}]
sudo pkill -f "snort -c"
sudo pkill -f "tcpdump -i eth0"
sudo pkill -f "testbed_inference.py"
\end{lstlisting}

\textbf{Theo dõi log:}

\begin{lstlisting}[language=bash, caption={Kiểm tra log Snort và ML alerts real-time}]
# Snort alerts
tail -f /var/log/snort/alert_fast.txt

# ML alerts
tail -f /var/log/ids/ml_alerts.log
\end{lstlisting}

\textbf{Yêu cầu phần cứng tối thiểu:} 8 GB RAM (CICFlowMeter Java heap + PyTorch), CPU Intel Core i5 Gen 8 trở lên. GPU không bắt buộc cho inference --- CPU inference đủ với batch nhỏ (dưới 1 giây mỗi batch 128 flows).

\end{document}
